% =========================================================================
% Secao 2 - Trabalhos previos e a divulgacao do INEP
% =========================================================================

\section{Trabalhos prévios e divulgação oficial}
\label{sec:tres_eras}

A construção de indicadores de fluxo escolar a partir de pesquisa
domiciliar não é nova no Brasil. O esforço mais conhecido nesse
sentido é o do grupo PROFLUXO, proposto por
\cite{Fletcher1985_modelo} e desenvolvido na Fundação Cesgranrio por
Sergio Costa Ribeiro, Ruben Klein e colaboradores. A intuição central
do grupo era que, em um único corte de uma pesquisa domiciliar como
a PNAD, a distribuição conjunta de idade e série de todos os
estudantes do país permite resolver um sistema de equações de fluxo
de estado estacionário e recuperar as taxas de promoção, repetência
e evasão que produziriam aquela distribuição observada
\citep{FletcherRibeiro1989_profluxo}. O método e suas extensões
foram a base de uma série de estudos sobre o fluxo escolar
brasileiro, incluindo \cite{KleinRibeiro1991_censo},
\cite{Fletcher1998_eficaz} e a síntese de \cite{Klein2006_educacao},
e tiveram impacto duradouro nas estatísticas educacionais
brasileiras. O achado emblemático de \cite{Ribeiro1991_pedagogia}
sobre a pedagogia da repetência, em particular, mostrou que a taxa
real de reprovação era ordens de magnitude maior do que a reportada
pelo Censo Escolar, contribuição que orientou políticas educacionais
durante mais de uma década. Esta nota não é uma extensão do
PROFLUXO. Não há sobreposição entre os autores e os instrumentos
metodológicos são distintos. Onde o PROFLUXO usava um modelo de
estado estacionário para inferir fluxos não observáveis a partir de
um único corte, esta nota observa diretamente as transições entre
dois pontos no tempo no painel longitudinal da PNADC. A menção ao
PROFLUXO se justifica pelo lugar do grupo na história das
estatísticas educacionais brasileiras e pela similaridade do
projeto subjacente, ou seja, aproveitar a pesquisa domiciliar como
fonte alternativa ao registro administrativo, mas não pelo arcabouço
técnico, que é completamente diferente.

No campo da avaliação educacional, o grupo da UFMG liderado por
José Francisco Soares e Maria Teresa Gonzaga Alves desenvolveu uma
linha de pesquisa sobre trajetórias educacionais irregulares e
contexto escolar. \citet{SoaresAlvesFonseca2021_trajetorias} propõem
o uso de trajetórias educacionais como evidência da qualidade da
educação básica brasileira, articulando rendimento e fluxo em uma
medida unificada, ao passo que \citet{AlvesSoares2013_contexto}
discutem como o contexto escolar gera condições desiguais para a
efetivação de políticas de avaliação, com implicações diretas para
a interpretação de indicadores de fluxo. Esses trabalhos
complementam a literatura PROFLUXO ao introduzir uma dimensão
escola-específica, ao passo que \citet{RianiRiosNeto2008_background}
mostram a importância do background familiar em relação ao perfil
escolar do município para o resultado educacional brasileiro. Outro
esforço aplicado importante é o de \citet{Fernandes2011_atratividade},
que usou dados mensais da Pesquisa Mensal de Emprego (PME) do IBGE
para mostrar como o número de matrículas evolui mês a mês ao longo
do ano letivo, separando alunos com e sem defasagem idade-série. As
figuras de Fernandes mostram que, para alunos em dia, a evasão
ocorre quase exclusivamente entre séries; para alunos defasados, há
queda contínua ao longo do ano letivo, indicando abandono intra-ano.
Esse padrão informa diretamente a discussão da
Seção~\ref{sec:agregados} desta nota e é replicado na
Figura~\ref{fig:cohort_fluxo}, agora com dados trimestrais da PNADC
e cobertura nacional, ao passo que a PME cobria apenas seis regiões
metropolitanas e foi descontinuada em 2016.

A literatura internacional fornece referências úteis para a
construção e a interpretação de indicadores de fluxo. Nos Estados
Unidos, o \emph{National Center for Education Statistics} produz
sistematicamente indicadores de \emph{dropout} e de
\emph{graduation rate} apurados a partir da combinação de pesquisas
domiciliares (Current Population Survey) e registros
administrativos. A Organização para a Cooperação e Desenvolvimento
Econômico (OCDE), em sua publicação anual \emph{Education at a
Glance}, sintetiza indicadores comparáveis para os países membros
e parceiros, incluindo taxas de conclusão por nível educacional e
indicadores de defasagem. A construção brasileira proposta nesta
nota dialoga com esses padrões internacionais ao usar pesquisa
domiciliar (PNADC) como fonte complementar ao registro
administrativo (Censo Escolar), com a vantagem de permitir
desagregação socioeconômica direta, propriedade que tipicamente
falta em registros administrativos.

Do lado administrativo, o INEP \citep{INEP2017_fluxo} introduziu o
CPF como identificador único do aluno no Censo Escolar em 2007 e
descreveu uma metodologia de acompanhamento longitudinal individual
explorando esse identificador estável entre Censos consecutivos. A
partir disso, definiu quatro indicadores de fluxo entre dois anos
consecutivos $t$ e $t+1$, a saber, promoção (matrícula em série
superior), repetência (matrícula na mesma série), migração para
EJA e evasão (sem matrícula em $t+1$). Por construção, as quatro
taxas somam um. A publicação dessas taxas, contudo, tem sido
errática. A nota técnica original \citep{INEP_transicao} cobriu o
período 2007--2016, mostrando, por exemplo, queda da evasão no
ensino médio de 11,1\% na transição 2013/2014 para 9,1\% na
transição 2016/2017. A série foi posteriormente estendida em duas
rodadas tardias até a transição 2021--2022, publicada em julho de
2025. Em junho de 2026, contudo, não há transições divulgadas para
2022--2023 nem 2023--2024, em contraste forte com as taxas de
rendimento, que têm publicação anual praticamente ininterrupta.

A divulgação dos diferentes conjuntos de indicadores oficiais segue
ritmos distintos. Os Indicadores de Rendimento, que englobam
aprovação, reprovação e abandono, tiveram sua última divulgação em
agosto de 2025, referente ao ano letivo de 2024, com defasagem
típica de oito a dez meses. Os Indicadores de Trajetória, que cobrem
promoção, repetência, migração para EJA e evasão, tiveram sua
última divulgação em julho de 2025, referente à transição
2021--2022, e a defasagem desse indicador ao presente, em junho de
2026, é portanto de quatro anos. O ritmo é errático, com pausas
longas entre rodadas e sem calendário público de atualização. O
IDEB, que combina rendimento e SAEB, foi divulgado em 2024 com
referência a 2023, defasagem típica de doze meses. A PNADC
trimestral, em comparação, tem divulgação em curso com defasagem
máxima de quatro a cinco meses em relação ao trimestre de
referência. A descontinuidade dos Indicadores de Trajetória é o
principal vácuo de informação que esta nota busca preencher.

A descontinuidade na divulgação tem consequências práticas. Em
primeiro lugar, gestores e pesquisadores que, em 2024--2026,
precisam monitorar evasão, repetência e promoção, dispõem apenas
das séries do Censo Escolar de rendimento, que captam abandono mas
não distinguem repetência efetiva da promoção. Em segundo lugar, a
defasagem de mais de uma década entre o último dado oficial de
trajetória (2016/2017) e o presente impossibilita avaliações
longitudinais de políticas como o Programa Bolsa Família, o
Auxílio Brasil e, mais recentemente, o Pé-de-Meia
(Lei 14.818/2024). Em terceiro lugar, a heterogeneidade
socioeconômica do fluxo escolar permanece invisível nos dados
oficiais publicados, ainda que essa heterogeneidade seja
precisamente o que motiva as políticas focalizadas em jovens em
situação de pobreza.

As transições 2019--2020, 2020--2021 e 2021--2022 cobrem o período
pandêmico, quando o conceito operacional de matrícula ativa e
frequência escolar sofreu rupturas profundas. As escolas brasileiras
tiveram aulas presenciais suspensas em 2020 com retorno gradual e
heterogêneo em 2021. Muitos sistemas adotaram aprovação automática
ou critérios excepcionais de progressão, distorcendo a interpretação
dos indicadores de fluxo. As taxas oficiais do INEP nesses anos
refletem essas convenções administrativas mais do que o processo
regular de aprendizagem e progressão. Adicionalmente, o IBGE
conduziu parte das entrevistas da PNADC por telefone, reduzindo o
tamanho amostral efetivo em até quarenta por cento em alguns
trimestres. Por essas razões, nas seções subsequentes destacamos
visualmente o período 2020--2021 nas figuras e excluímos esses anos
da média multianual reportada.

Do lado da pesquisa domiciliar, a publicação dos identificadores
longitudinais individuais da PNADC pelo IBGE, a partir de 2019,
permite seguir o mesmo indivíduo em até cinco trimestres
consecutivos, perfazendo quinze meses no total. Cada domicílio
entrevistado entra em uma fileira de rotação, o painel, e é
visitado em cinco trimestres antes de sair da amostra. Combinando o
identificador domiciliar estável (UF, UPA, número de seleção,
número do painel) com validações de consistência individual (sexo,
idade e raça), é possível ligar registros de uma mesma pessoa
entre as cinco visitas com taxa de acerto superior a noventa por
cento em pré-testes. Esta nota explora exatamente essa
possibilidade, observando diretamente, sem modelo de estado
estacionário, as transições entre anos consecutivos com janela de
medição restrita aos trimestres $Q_2$ e $Q_3$, e as transições
intra-ano dentro do mesmo ano calendário, estimando a partir delas
os indicadores de fluxo. A próxima seção descreve em detalhe a
construção do painel e as definições operacionais.
